%% HAND-WRITTEN circuitikz figure -- NOT generated by maestro2tex. The tool only places
%% it, via the "order" list in configs/tb_anatop_tia_rprog_corners.json; the `sch_`
%% prefix keeps it clear of the generated fig_*/tab_* namespace.
%%
%% Depicts: the segment/bypass structure of castalia/anatop_tia_rprog -- the series
%% chain A -> 3-unit base -> six binary segments (32/16/8/4/2/1 units, MSB first) ->
%% four 64-unit thermometer segments -> B, each switched segment shunted by a CMOS
%% transmission gate with a local inverter (anatop_tia_rprog_sw, pass pair multi=4,
%% for the binary segments; anatop_tia_rprog_sw8, multi=8, for the thermometer
%% segments). The switch-cell internals (NMOS+PMOS pass pair, local inverter
%% EnBar->En) are stated in the caption, not drawn.
%%
%% Topology transcribed on 2026-08-21 from the executed spec-run netlist
%%   simruns/tb_anatop_tia_rprog_specrun/castalia/tb_anatop_tia_rprog/maestro/
%%     results/maestro/Ocean.0/psf/rp_code/netlist/input.scs
%% (subckts anatop_tia_rprog / _sw / _sw8: R_BASE<2:0> A->RBaseOut; R_B5..R_B0
%% RBaseOut->BinOut with TG_B5..TG_B0 in parallel per segment; R_TH0..R_TH3
%% BinOut->B with TG_T0..TG_T3; R_DUMMY<2:0> tied to vss. Switch cells: M0 nch
%% pass gate=En, M1 pch pass gate=EnBar, M2/M3 local inverter EnBar->En) and
%% cross-checked against the builder skill/castalia_tia_rprog_build.il.
%% Per the section's rule, no device geometry is stated -- only multiplicities.
\begin{figure}[htbp]
  \centering
  \resizebox{\linewidth}{!}{%
\ctikzset{bipoles/length=0.85cm, resistors/scale=0.85, transistors/scale=1.0, logic ports/scale=0.6}
\begin{circuitikz}[
    every node/.append style={font=\footnotesize},
    nlab/.style={font=\scriptsize, inner sep=1.5pt, fill=white},
    ctlab/.style={font=\scriptsize, inner sep=1.5pt},
    devlab/.style={font=\scriptsize, text=black!60, inner sep=1pt},
    railnode/.style={circle, fill=black, inner sep=1.05pt},
    grp/.style={draw=black!45, dashed, line width=0.6pt, rounded corners=4pt},
  ]

% ===== upper row: A -> base -> binary segments ============================
\def\segw{2.55}
\draw (0,0) node[railnode]{} node[nlab, above=3pt]{\texttt{A}}
      to[R, l=$3R_u$] (2.0,0) node[railnode]{};
\node[devlab, below] at (1.0,-0.32) {base};

\foreach \i/\nu/\ctl in {0/32/5, 1/16/4, 2/8/3, 3/4/2, 4/2/1, 5/1/0} {
  % series segment
  \draw ({2.0+\i*\segw},0) to[R, l=$\nu R_u$] ({2.0+(\i+1)*\segw},0)
        node[railnode]{};
  % bypass transmission gate, drawn as a switch; cell detail at right
  \draw ({2.0+\i*\segw},0) -- ({2.0+\i*\segw},-1.15)
        to[nos] ({2.0+(\i+1)*\segw},-1.15) -- ({2.0+(\i+1)*\segw},0);
  \node[ctlab, below] at ({2.0+(\i+0.5)*\segw},-1.45) {\texttt{ResEn<\ctl>}};
}
\draw[grp] (1.7,0.75) rectangle (17.6,-2.1);
\node[devlab, anchor=south west] at (1.75,0.82)
      {binary segments --- bypass \texttt{anatop\_tia\_rprog\_sw} (pass pair $\times 4$)};

% ===== lower row: thermometer segments -> B ===============================
% bridge from BinOut (end of the binary chain) down to the thermometer chain
\draw (17.3,0) -- (18.05,0) -- (18.05,-3.6) -- (17.3,-3.6);
\node[railnode] at (17.3,-3.6){};
\foreach \i/\ctl in {0/0, 1/1, 2/2, 3/3} {
  \draw ({17.3-\i*4.1},-3.6) to[R, l_=$64R_u$] ({17.3-(\i+1)*4.1},-3.6)
        node[railnode]{};
  \draw ({17.3-\i*4.1},-3.6) -- ({17.3-\i*4.1},-4.75)
        to[nos] ({17.3-(\i+1)*4.1},-4.75) -- ({17.3-(\i+1)*4.1},-3.6);
  \node[ctlab, below] at ({17.3-(\i+0.5)*4.1},-5.05) {\texttt{ThEn<\ctl>}};
}
\draw (0.9,-3.6) -- (0,-3.6) node[railnode]{} node[nlab, above=3pt]{\texttt{B}};
\draw[grp] (0.65,-2.85) rectangle (17.6,-5.7);
\node[devlab, anchor=south west] at (0.7,-2.78)
      {thermometer segments --- bypass \texttt{anatop\_tia\_rprog\_sw8} (pass pair $\times 8$)};

\end{circuitikz}
  }
  \caption[{Segment and bypass-switch structure of the programmable feedback
  resistor \texttt{anatop\_tia\_rprog}.}]{Segment and bypass-switch structure of
  \texttt{anatop\_tia\_rprog}, transcribed from the executed spec-run netlist.
  The unit \(R_u\) is one series poly resistor; from terminal \texttt{A} the
  chain runs through the never-switched three-unit base, the six binary-weighted
  segments of 32 down to 1 unit, and the four identical 64-unit thermometer
  segments to terminal \texttt{B}. Each switched segment is shunted by a CMOS
  transmission gate with a local inverter: the control pin \texttt{EnBar} is
  the \texttt{ResEn}/\texttt{ThEn} bit itself and drives the PMOS pass gate
  directly and the NMOS pass gate through the inverter, so a bit at 1 turns
  both pass devices off and inserts the segment. The pass pair has
  multiplicity 4 in the binary switch (\texttt{\_sw}) and 8 in the thermometer
  switch (\texttt{\_sw8}), which halves the thermometer gates' on-resistance
  where up to six binary gates close at a seam. Three grounded dummy units
  accompany the 322 active units for matching.}
  \label{fig:tiarprog-switch}
\end{figure}
